\documentclass[a4paper]{article}

% Español (México) con XeLaTeX: fontspec para fuentes Unicode, polyglossia
% para la división silábica y los nombres del documento en la variante mexicana.
\usepackage{fontspec}
\usepackage{polyglossia}
\setdefaultlanguage[variant=mexican]{spanish}
% Los nombres chinos van como «pinyin (original)»: xeCJK da a los caracteres
% Han su propia fuente; sin ella XeLaTeX los omite con un aviso.
% FandolSong no cubre algunos caracteres de nombres (U+73FA); WenQuanYi Zen Hei sí.
\usepackage[AutoFallBack=true]{xeCJK}
\setCJKmainfont{FandolSong-Regular.otf}
\setCJKfallbackfamilyfont{\CJKrmdefault}{WenQuanYi Zen Hei}
% polyglossia no traduce los nombres de listings.
\AtBeginDocument{\providecommand{\lstlistingname}{}\renewcommand{\lstlistingname}{Listado}%
  \providecommand{\lstlistlistingname}{}\renewcommand{\lstlistlistingname}{Índice de listados}}
\usepackage{amsmath,amssymb}
\usepackage{graphicx}
\usepackage[margin=2.5cm]{geometry}
\usepackage[most]{tcolorbox}
\usepackage{etoolbox}
\usepackage{subcaption}
\usepackage{listings}
\usepackage{hyperref}
\usepackage{booktabs}
\usepackage{float}

\newtcolorbox{knowledgebox}[1]{
    enhanced, colback=blue!5!white, colframe=blue!75!black, colbacktitle=blue!75!black,
    coltitle=white, fonttitle=\bfseries, title={#1}, boxrule=1pt, sharp corners
}

\newtcolorbox{importantbox}[1]{
    enhanced, colback=yellow!10!white, colframe=yellow!80!black, colbacktitle=yellow!80!black,
    coltitle=black, fonttitle=\bfseries, title={#1}, sharp corners
}

\newtcolorbox{warningbox}[1]{
    enhanced, colback=red!5!white, colframe=red!75!black, colbacktitle=red!75!black,
    coltitle=white, fonttitle=\bfseries, title={#1}, sharp corners
}

\lstset{
    basicstyle=\ttfamily\small,
    keywordstyle=\color{blue},
    stringstyle=\color{red!60!black},
    commentstyle=\color{green!60!black},
    breaklines=true,
    frame=single,
    numbers=left,
    numberstyle=\tiny\color{gray},
    captionpos=b,
    extendedchars=true
}

\newcommand{\notetitle}{CS224R Lecture 18: Fronteras del RL profundo y métodos de investigación}
\newcommand{\noteauthors}{Elaboradas a partir de materiales públicos del curso}
\newcommand{\notedate}{Primavera de 2025}
\newcommand{\videochannel}{Stanford Online}
\newcommand{\videopublishdate}{2025}
\newcommand{\videoduration}{Aproximadamente 80 minutos}
\newcommand{\videourl}{https://cs224r.stanford.edu/}
\newcommand{\videocoverpath}{cover.jpg}
\newcommand{\repourl}{https://github.com/hqhq1025/ai-course-notes}


% Footer with repo info
\usepackage{fancyhdr}
\pagestyle{fancy}
\fancyhf{}
\fancyfoot[C]{\thepage}
\fancyfoot[R]{\footnotesize\color{gray}\href{\repourl}{AI Course Notes}}
\renewcommand{\headrulewidth}{0pt}
\renewcommand{\footrulewidth}{0pt}

\begin{document}
\begin{titlepage}
\centering
{\Large Notas del curso\par}
\vspace{1.2cm}
{\huge\bfseries \notetitle\par}
\vspace{0.8cm}
{\large \noteauthors\par}
\vspace{0.3cm}
{\large \notedate\par}
\vspace{1.1cm}

\ifdefempty{\videocoverpath}{
{\small No se proporcionó la imagen de portada.\par}
}{
\includegraphics[width=0.82\textwidth,height=0.45\textheight,keepaspectratio]{\videocoverpath}\par
}

\vfill
\begin{tcolorbox}[width=0.9\textwidth, colback=black!2!white, colframe=black!60, sharp corners]
\textbf{Autor o canal del video}: \videochannel\par
\textbf{Fecha de publicación}: \videopublishdate\par
\textbf{Duración del video}: \videoduration\par
\textbf{Ponente}: Chelsea Finn\par
\textbf{Página del curso}: \href{https://cs224r.stanford.edu/}{cs224r.stanford.edu}\par
\textbf{Página del proyecto}: \href{\repourl}{\nolinkurl{\repourl}}\par
\end{tcolorbox}
\end{titlepage}

\tableofcontents
\newpage

%% --- Inicio del contenido --- %%

\section{Parte 1: problemas abiertos y fronteras del RL profundo}

En esta clase, Chelsea Finn resume los problemas abiertos y las líneas de frontera del RL profundo, y ofrece recomendaciones prácticas para investigar en RL profundo.

\subsection{Desafíos en la definición del problema}

\subsubsection{Diseño de la recompensa}

\begin{importantbox}{La dificultad de fondo del diseño de la recompensa}
En las aplicaciones reales, diseñar la función de recompensa correcta suele ser más difícil que resolver el propio problema de RL:
\begin{itemize}
    \item \textbf{Recompensa dispersa}: en la mayoría de los pasos de tiempo no hay retroalimentación significativa
    \item \textbf{Recompensa mal especificada}: la recompensa diseñada puede producir conductas no previstas (reward hacking)
    \item \textbf{Conflicto entre objetivos múltiples}: es difícil ajustar los pesos entre varios objetivos de recompensa
\end{itemize}
\end{importantbox}

\begin{warningbox}{Reward Hacking}
Cuando la función de recompensa es imperfecta, el agente de RL encuentra formas «legítimas pero no previstas» de maximizar la recompensa. Casos clásicos:
\begin{itemize}
    \item En un juego de carreras, recolectar monedas de recompensa en lugar de completar la pista
    \item Un robot que maximiza la recompensa de velocidad vibrando en lugar de caminar
    \item Un LLM que maximiza la recompensa por longitud repitiendo su salida
\end{itemize}
El reward hacking es uno de los desafíos centrales del problema de alineación (alignment) en RL.
\end{warningbox}

\begin{figure}[H]
\centering
\includegraphics[width=0.9\textwidth]{slides-images/slide-05.jpg}
\caption{El curso ilustra con chatbots, robots que doblan ropa y sistemas de recomendación que, en cuanto una tarea carece de una recompensa confiable o de una métrica verificable objetivamente, el propio planteamiento del problema de RL se vuelve un tema de investigación.}
\end{figure}

\subsubsection{Diseño de los espacios de estados y de acciones}

Elegir una representación de estados y un espacio de acciones adecuados también es una cuestión de diseño importante:
\begin{itemize}
    \item El estado debe contener información suficiente (cumplir la propiedad de Markov), pero sin tener una dimensión demasiado alta
    \item La granularidad del espacio de acciones afecta la eficiencia del aprendizaje
    \item En los LLM, si el espacio de acciones es a nivel de token o a nivel de sentence
\end{itemize}

\subsubsection{Diseño de entornos y puntos de referencia}

\begin{knowledgebox}{Criterios de un buen punto de referencia de RL}
Un punto de referencia de RL útil debe:
\begin{itemize}
    \item Capturar los desafíos centrales de las aplicaciones reales
    \item Permitir iterar con rapidez (con un costo computacional moderado)
    \item Tener métricas de evaluación significativas y confiables
    \item Correlacionarse con el desempeño en el mundo real
\end{itemize}
El desafío actual es que muchos puntos de referencia populares de RL están desconectados de las necesidades de las aplicaciones reales.
\end{knowledgebox}

\subsection{Desafíos en los métodos}

\subsubsection{Eficiencia de muestreo}

Pese a las numerosas mejoras, la eficiencia de muestreo del RL profundo sigue siendo muy inferior a la del aprendizaje humano.

\subsubsection{Capacidad de generalización}

\begin{importantbox}{El problema de la generalización en RL}
La capacidad de generalización de los métodos de RL actuales es limitada:
\begin{itemize}
    \item El desempeño cae abruptamente en entornos nuevos fuera del entorno de entrenamiento
    \item Son excesivamente sensibles a cambios mínimos en la función de recompensa
    \item Les cuesta transferir lo aprendido en una tarea a tareas relacionadas
\end{itemize}
Mejorar la capacidad de generalización del RL es indispensable para desplegarlo en el mundo real.
\end{importantbox}

\subsubsection{Sistemas multiagente y juegos}

El RL multiagente enfrenta desafíos adicionales:
\begin{itemize}
    \item El entorno no es estacionario (las políticas de los demás agentes cambian)
    \item El equilibrio entre coordinación y competencia
    \item El aprendizaje de protocolos de comunicación
\end{itemize}

\begin{figure}[H]
\centering
\includegraphics[width=0.9\textwidth]{slides-images/slide-12.jpg}
\caption{La ponente sitúa la video generation / world model entre las fronteras de los métodos: aportan priors ricos, pero cuando se usan fuera de la distribución de la política se distorsionan por errores físicos y por acciones fuera de distribución.}
\end{figure}

\subsection{Desafíos en el despliegue y la evaluación}

\subsubsection{Seguridad y confiabilidad}

\begin{warningbox}{Riesgos de desplegar sistemas de RL}
Desplegar sistemas de RL en el mundo real implica múltiples riesgos:
\begin{itemize}
    \item Los estados fuera de distribución (out-of-distribution) pueden provocar conductas catastróficas
    \item El proceso de decisión de una política de RL es difícil de explicar
    \item Los ataques adversarios pueden aprovechar las debilidades de la política
    \item Los efectos a largo plazo son difíciles de predecir
\end{itemize}
\end{warningbox}

\subsubsection{Reproducibilidad}

La reproducibilidad de los experimentos de RL profundo ha sido siempre motivo de preocupación:
\begin{itemize}
    \item Sensibilidad a la semilla aleatoria
    \item Alto costo del ajuste de hiperparámetros
    \item Las diferencias entre distintas implementaciones pueden ser muy grandes
\end{itemize}

\subsection{Resumen de la sección}

El RL profundo enfrenta desafíos en todos los niveles, desde la definición del problema hasta los métodos y el despliegue. El diseño de la recompensa, la capacidad de generalización y la seguridad son los problemas abiertos más importantes.
\section{Parte 2: cómo hacer investigación en RL profundo}

\subsection{Actitud y metodología de investigación}

\begin{importantbox}{Recomendaciones de investigación de Chelsea Finn}
\begin{enumerate}
    \item \textbf{Partir del problema}: primero entender qué problema se quiere resolver y después pensar qué método usar. No «tener un martillo y buscar clavos».
    \item \textbf{Entender antes de innovar}: comprender a fondo las ventajas y desventajas de los métodos existentes; la innovación surge de forma natural.
    \item \textbf{Lo simple primero}: empezar por el método o el experimento más sencillo y aumentar la complejidad de manera gradual.
    \item \textbf{Construir intuición}: usar entornos sencillos para formar intuición sobre el comportamiento del algoritmo y luego generalizarla a escenarios complejos.
\end{enumerate}
\end{importantbox}

\subsection{Diseño experimental}

\begin{knowledgebox}{Buenas prácticas para experimentos de RL}
\begin{itemize}
    \item \textbf{Control de variables}: cambiar un solo factor a la vez.
    \item \textbf{Ejecuciones con varias semillas}: al menos de 3 a 5 semillas aleatorias.
    \item \textbf{Líneas base significativas}: elegir líneas base sólidas y no hombres de paja.
    \item \textbf{Experimentos de ablación}: verificar la contribución de cada decisión de diseño.
    \item \textbf{Herramientas de diagnóstico}: monitorear las curvas de aprendizaje, el error de estimación de la función de valor, la entropía de la política, etc.
\end{itemize}
\end{knowledgebox}

\begin{figure}[H]
\centering
\includegraphics[width=0.9\textwidth]{slides-images/slide-15.jpg}
\caption{El curso concreta «cómo escalar RL» en tres preguntas de diseño experimental: por lotes contra en línea, precisión de la función de valor y frecuencia de actualización. Estas preguntas determinan si un algoritmo puede escalar realmente a escenarios de horizonte largo y de modelos grandes.}
\end{figure}

\subsection{Depuración de algoritmos de RL}

Depurar algoritmos de RL es mucho más difícil que en el aprendizaje supervisado. Estrategias de depuración que propone la ponente:

\begin{enumerate}
    \item \textbf{Validar en entornos sencillos}: confirmar primero que el algoritmo funciona correctamente en CartPole/Pendulum.
    \item \textbf{Revisar la curva de recompensa}: la recompensa debe crecer de forma monótona (o aproximadamente monótona).
    \item \textbf{Revisar la función de valor}: la estimación de la función de valor debe acercarse al retorno real.
    \item \textbf{Visualizar el comportamiento de la política}: observar qué hace realmente el agente.
    \item \textbf{Revisar los gradientes}: asegurarse de que los gradientes no exploten ni se desvanezcan.
\end{enumerate}

\begin{warningbox}{Errores comunes al depurar RL}
\begin{itemize}
    \item Una scale inadecuada de la recompensa provoca que la función de valor diverja.
    \item Una elección inadecuada del factor de descuento $\gamma$.
    \item Omitir la normalización o estandarización de los datos.
    \item Una exploration insuficiente hace que la política converja a un óptimo local.
    \item La aleatoriedad puede ocultar un bug: aun con un bug, a veces RL logra aprender una política «aceptable».
\end{itemize}
\end{warningbox}

\subsection{Cómo elegir una línea de investigación}

La ponente ofrece recomendaciones para elegir una línea de investigación:
\begin{itemize}
    \item Enfocarse en \textbf{problemas realmente importantes} y no en temas de moda.
    \item Buscar la \textbf{brecha} entre la teoría y la práctica.
    \item Pensar entre disciplinas (NLP $\times$ RL, robotics $\times$ RL).
    \item Conversar con profesionales del área para conocer las necesidades reales.
\end{itemize}

\subsection{Resumen de la sección}

Una buena investigación en RL requiere una definición clara del problema, un diseño experimental riguroso y estrategias de depuración eficaces. Lo simple primero y el control de variables son la metodología central.
\section{Síntesis y ampliación}

\begin{enumerate}
    \item Los problemas abiertos del RL profundo abarcan tres niveles: la definición del problema, los métodos y el despliegue
    \item El diseño de recompensas (para evitar el reward hacking), la capacidad de generalización y la seguridad son las líneas más importantes
    \item La buena investigación parte de entender el problema, no del método
    \item Los experimentos de RL requieren varias semillas aleatorias, líneas base sólidas y experimentos de ablación
    \item La depuración de RL requiere un método sistemático; la visualización y la verificación en entornos simples son esenciales
    \item La integración entre áreas (RL + LLM, RL + robotics) es la frontera más activa en la actualidad
\end{enumerate}

\begin{knowledgebox}{Lecturas adicionales}
\begin{itemize}
    \item Henderson et al., ``Deep Reinforcement Learning that Matters,'' AAAI 2018
    \item Amodei et al., ``Concrete Problems in AI Safety,'' arXiv 2016
    \item Engstrom et al., ``Implementation Matters in Deep Policy Gradients,'' ICLR 2020
    \item Andrychowicz et al., ``What Matters in On-Policy Reinforcement Learning? A Large-Scale Empirical Study,'' ICLR 2021
\end{itemize}
\end{knowledgebox}

\end{document}
